\documentclass[
  paper=a4,
  fontsize=12pt,
  DIV=16,
  headheight=52pt,
  footheight=45pt,
  headinclude,
  parskip=full,
]{scrartcl}

\usepackage{fontspec}
\setsansfont{Source Sans Pro}
\renewcommand{\familydefault}{\sfdefault}


\usepackage{polyglossia}
\setmainlanguage{german}

\usepackage{csquotes}

\usepackage{ulem}
\usepackage{microtype}
\usepackage{graphicx}

\usepackage{scrlayer-scrpage}
\pagestyle{scrheadings}
\setkomafont{pagehead}{\normalfont}
\setkomafont{pagefoot}{\normalfont\footnotesize}
\ohead{\includegraphics[height=1.5cm]{../../../logos/pep.pdf}}
\ihead{%
  \large\bfseries PeP et al.\ e.\,V.\\%
  Begabtenförderungsprogramm –\\
  Going Abroad - Purdue%
}
\cfoot{}
\ifoot{%
  PeP et al.~e.\,V.\\
  Begabtenförderung\\
  \url{www.pep-dortmund.org}
}

\usepackage[hyphens]{url}

\usepackage{titling}
\usepackage{booktabs}

\usepackage{xcolor}
\usepackage{calc}
\usepackage[colorlinks=true,urlcolor=blue!50!black]{hyperref}

\date{\today}

\begin{document}
\textbf{\Huge\sffamily Bachelorarbeit an der Purdue University}\\[0.5\baselineskip]
\textbf{\Large\sffamily Eine Anleitung was alles im Vorfeld tun ist}

Du hast die Zusage zur Teilnahme am Going Abroad Programm für Purdue erhalten.
Mit der alleinigen Zusage ist es leider nicht getan.
Es gibt noch einige wichtige Aufgaben, die du vorher erledigen musst.
Im folgenden werden dir die wichtigsten Fragen beantwortet.
Solltest du nicht ausreichend informiert sein, oder noch weitere Fragen haben,
kannst du dich an Marie Schmitz (CP-03-188) persönlich oder per Mail
(\href{mailto:marie.schmitz@tu-dortmund.de}{marie.schmitz@tu-dortmund.de}) wenden.
Alternative kannst du auch dem ehmaligen Stipendiat Steven Becker eine Mail
schreiben (\href{mailto:steven.becker@tu-dortmund.de}{steven.becker@tu-dortmund.de}).
Falls deine Fragen direkt mit Purdue zu tun haben, wie etwa der Unterkunft oder
dem ISS (\textit{Office of International Students and Scholars}), kannst du dich auch direkt an Andreas Jung (\href{mailto:anjung@purdue.edu}{anjung@purdue.edu}) wenden.
Wenn du nicht sicher bist, wer deine Fragen besser beantworten kann, schicke die Mail einfach an Marie, Andreas und Steven.


\subsection*{Kurzfassung}
Die folgenden Dinge müssen am Anfang recht schnell passieren, damit der oft langwierige Visumsantrag rechtzeitig abgeschlossen werden kann.
\begin{enumerate}
  \item \sout{\textbf{Schreibe eine Mail an Andreas Jung} (\href{mailto:anjung@purdue.edu}{anjung@purdue.edu}).}
    Er benötigt
    \begin{itemize}
      \item dein Geburtsdatum und den vollen Namen, wie er auch im Reisepass auftaucht
      \item deine permanente Postadresse (damit Vertragsunterlagen ankommen)
      \item sowie den Zeitraum des Programms
    \end{itemize}
    Das Programm ist standardmäßig für einen Zeitraum von 3 Monaten, vom 15.5.2018 – 15.8.2018 festgelegt.
    Bitte bestätige den Zeitraum oder handle mit Andreas eine Alternative aus.
  \item \textbf{Beantrage einen Reisepass im Eilverfahren}, falls du noch keinen besitzt.
    Der Reisepass ist unbedingt nötig für alle weiteren Schritte und sollte sofort beantragt werden.
    Sobald du den Reisepass vorliegen hast, brauchen wir \textbf{eine Kopie}.
  \item Außerdem muss die \textbf{Einverständniserklärung der Fakultät} vom \textbf{betreuenden Professor an der TU Dortmund unterschreiben lassen} und bei Henning im Büro abgeben.
  Diese Erklärung ist sehr wichtig und wird schnellstmöglich von uns benötigt.
  Erst mit ihr kann der Visumsprozess an der Purdue University starten, sowie die Unterkunft vor Ort organisiert werden.
  Zudem benötigen wir das ausgefüllte \textbf{Visiting Scholar Information Form}, dass du ebenfalls bei uns abgibst.
\end{enumerate}


\subsection*{Weitere Informationen}
\subsubsection*{Wie bekomme ich das richtige Visum für die USA?}
Das für dich richtige Visum ist das J1 Visum (\url{https://j1visa.state.gov}), das zu einen Austauschprogramm an einer US amerikanischen Universität berechtigt.
Der Prozess dieses zu erhalten ist etwas komplexer und wird dir im Folgenden dargelegt.

Zuerst musst du wie oben bereits beschrieben die Einverständniserklärung der Fakultät von deinem betreuenden Professor unterschreiben lassen und bei Marie im Büro
abgeben.
Diese Erklärung ist sehr wichtig und wird schnellstmöglich von uns benötigt.
Erst mit ihr können wir den Visumsprozess an der Purdue University starten, sowie deine Unterkunft vor Ort organisieren.
Das \textbf{Visiting Scholar Information Form} musst du mit den angefragten Daten füllen und ebenfalls bei uns abgeben.

Sobald wir die Dokumente unterschrieben von dir zurück bekommen haben, werden diese an Andreas Jung in die USA weitergeleitet.
Er wird dann an der Purdue Universität dafür sorgen, dass du dort auch angenommen wirst.
Dazu wird er bei verschiedenen Stellen, wie dem \textit{Business Office}, dem \textit{Visa and Immigration Office} und dem \textit{Office of International Students and Scholars (ISS)}, alles Notwendige in die Wege leiten.

Mit den unterschriebenen Dokumenten musst du uns auch eine \textbf{Kopie des Reisepasses} einreichen, damit auf amerikanischer Seite die Dokumente für das Visum mit den richtigen Daten ausgefüllt werden können.
Falls du deinen Reisepass noch beantragen musst, brauchen wir deinen \textbf{vollständigen Namen, wie er später im Pass steht}, sowie \textbf{dein Geburtsdatum}.
Den Reisepass solltest du dann dennoch idealerweise im Eilverfahren beantragen.
Bei dem Beantragen des Reisepasses kannst du deine Reisepassnummer erfragen und diese
uns anschließend schonmal mitteilen.

Während dieser Vorbereitungen wirst du ein paar Mails erhalten, die du
ebenfalls zügig beantworten solltest.
Unter anderem wirst du dazu aufgefordert werden die \textbf{MyISS J Request E-forms} auszufüllen.
Hier wirst du zu deinen bisherigen Reisen in die USA und in andere Länder gefragt, musst Daten zu deinen letzten Arbeitgebern und zum Studium, sowie deine Passnummer angeben.
Es wird auch nach deinen Reisedaten gefragt werden.
Hier reicht es erst einmal einen groben Zeitraum von Mai 2019 bis August 2019 anzugeben.

Letzlich wird es auch noch ein Interview in der amerikanischen Botschaft in Berlin oder Frankfurt geben, bei dem geprüft wird, ob du das Visum erhälst oder nicht.
Weitere Informationen dazu gibt es hier: \url{https://de.usembassy.gov/de/visa/}

Wenn du hier noch weitere Fragen hast, kannst du außerdem beim \textit{US Department of State} noch weitere Informationen erhalten: \url{https://travel.state.gov/content/visas/en/forms/ds-160--online-nonimmigrant-visa-application/frequently-asked-questions.html}
Natürlich stehen Marie, Andreas oder Steven dir ebenfalls helfend zu Seite.


\subsubsection*{Warum muss ich mich krankenversichern?}
Eine Sommerschule in den USA ist nur mit einem J-1 Visum möglich.
Der Visumsprozess erfordert zwingend das Vorhandensein einer \textbf{US-tauglichen Krankenversicherung}.
In den Staaten reicht deine deutsche Versicherung nicht aus und du musst selbst eine Zusatzversicherung abschließen.
Hier gibt es verschiedene Möglichkeiten:
\begin{itemize}
  \item Steven hat sich über \href{https://www.dr-walter.com/index.html}{Dr. Walter}
        versichert.
  \item Es kann auch eine Krankenversicherung in den USA zugekauft werden, hier
  gibt es zum Beispiel die folgenden Varianten:
  \begin{itemize}
    \item \url{www.compassstudenthealthinsurance.com/compare_international_insurance_plans.php}
    Eine US Krankenversicherung ist immer mit sehr viel Auswahl verbunden.
    Wir empfehlen eine most-inclusive insurance, um bei tatsächlicher Krankheit den hohen Kosten des US Gesundheitssystems vorzubeugen.
    Compass Benchmark Plus in der entsprechenden Altersklasse liegt - mit Selbstbeteiligung - bei etwa 200 € pro Monat.
    \item Das ISS von Purdue empfiehlt diese zwei Optionen:
    \begin{itemize}
      \item \url{www.compassstudenthealthinsurance.com/}
      \item \url{www.isoa.org/}
    \end{itemize}
  \end{itemize}
\end{itemize}

\subsubsection*{Okay, aber muss ich das alles selber zahlen?}
Nein, dazu hast du dich ja bei uns auf das Austauschprogramm beworben.
Du wirst seitens der Purdue Universität einen Arbeitsvertrag bekommen der den Großteil deiner Kosten vor Ort decken sollte.
Über das Geld was du erhältst kannst du frei verfügen.
Allerdings musst du davon die Krankenversicherung und deine Unterkunft bezahlen.

Das Geld für Flug, Reisepass und Visum, musst du vorstrecken.
Sobald du jedoch eine Rechnung dafür hast, wird PeP dir das Geld erstatten.
Die genauen Beträge, die PeP im Zuge dieses Stipendiums unterstützend zahlt, kannst du im Stipendienvertrag nachschauen.
Hebe daher bitte unbedingt die Rechnungen auf, damit wir das Geld erstatten können.

\subsubsection*{Wie komme ich nach West Lafayette, Indiana?}
Mit dem Flugzeug geht es nach entweder nach Indianapolis, Indiana oder zum
Chicago O'Hare International Airport.
Von da aus am besten weiter mit dem Bus nach West Lafayette.
Eine empfohlene Busverbindung erhältst du nachdem wir die Flugdaten abgeklärt
haben.
Den Flughafen für Hin- und Ruckflug kannst du dir aussuchen.
Hast du ein passendes Angebot gefunden, so buchst du eigenständig nach Absprache mit uns.
Das Geld für den Flug und den Bus wird dir nach Erhalt der Rechnung überwiesen.

Der Aufenthalt in Unterkünften der Uni ist vom 15. Mai bis zum 15. August
möglich. Das ist genau der Zeitraum des Summer Break. Wenn du vor oder nach dem
Programm noch in den USA auf eigene Kosten reisen willst, wende dich noch einmal an uns, da lässt sich evtl. eine individuelle Regelung finden.


\subsubsection*{Wo schlafe ich?}
Während deines Aufenthaltes an der Purdue University wirst du in einem Einzelzimmer
in der Hawkins Hall untergebracht.
Hier kannst du deine Unterkunft schon einmal besichtigen und erste Fragen beantwortet bekommen: \url{https://www.housing.purdue.edu/Housing/Residences/Hawkins/index.html}
%Hat man einen Safe, also wie sieht es mit Wertsachen aus?

\subsubsection*{Wie werde ich versorgt?}
In deiner Unterkunft gibt es eine große Gemeinschaftsküche.
Dort kannst du dich mit anderen Studenten arrangieren und ihr könnt gemeinsam kochen.

Wenn du selbst nicht gerne kochst, ist in deiner Unterkunft ebenfalls eine
Dining Card enthalten.
Mit dieser kannst du die verschiedenen Dining Courts, Markets, Restaurants und Cafes am Campus besuchen.
Hier wirst du sicher deinen Favoriten finden.
Einen ersten Überblick kannst du dir hier verschaffen \url{https://dining.purdue.edu/AboutUs/map.html}.
Beliebte Orte findest du hier: \url{https://lafayettedaily.com/purdue-food/}.

\subsubsection*{Gibt es Möglichkeiten sich sportlich am Campus zu betätigen?}
An der Purdue University gibt es ein breites Sportangebot.
Informieren kannst du dich hier: \url{https://www.purdue.edu/recwell/index.php}.
Eine Mitgliedschaft musst du allerdings selbst aus dem Geld zahlen, welches du aus dem Arbeitsvertrag bekommst.

\subsubsection*{Wie läuft das mit der Bachelorarbeit?}
Dieser Teil ist mit der wichtigste, denn du möchtest ja auch vor allem deinen ersten akademischen Abschluss machen.
Du wirst deine Bachelorarbeit während deines Aufenthaltes an der Purdue University unter der Aufsicht und Betreuung von Andreas Jung schreiben.
Hier kannst du dir aussuchen ob du die Arbeit auf Deutsch oder auf Englisch verfassen willst.

Auch auf deutscher Seite hast du einen Professor, der dich betreuen wird.
Da er allerdings nicht vor Ort ist, kannst du ihm nur per Mail Fragen stellen.
Deine Arbeit musst du hinterher bei deinem Betreuer in Deutschland abgeben, denn du bist ja offiziell an der TU Dortmund eingeschrieben und machst auch dort deinen Abschluss.

Wenn du Fragen hast, die du weder Andreas Jung, noch deinem betreuenden Professor stellen willst, kannst du dich auch vertrauensvoll an PeP wenden.
Wir werden dir dann einen Ansprechpartner aus unseren Reihen zu Seite stellen,
der dir bei deinen Fragen weiterhelfen wird.

Deinen Bachelorvortrag musst du anschließend an deinen Auslandsaufenthalt ganz normal in Deutschland am Lehrstuhl deines deutschen Professors halten.
Du solltest dort - wie bei jedem anderen Bachelorvortrag auch - dein Thema
darstellen, sowie die Ergebnisse deiner Arbeit präsentieren.
Neben dem fachlichen Teil ist es durch die Teilnahme am Going Abroad Programm ebenfalls erforderlich, dass du etwa zehn bis fünfzehn Minuten über deine Erfahrungen und Arbeit an der Purdue University berichtest.

Die Benotung deiner Arbeit folgt den gleichen Regeln, als wenn du die Arbeit an der TU geschrieben hättest.
Der einzige Unterschied ist, dass Andreas Jung deinem betreuenden Professor sehr wahrscheinlich eine Empfehlung für die Note aussprechen wird.

\subsubsection*{Wie sieht der Arbeitsalltag aus?}
Um sich den Arbeitsalltag besser vorstellen zu können, sind hier zwei Berichte
von Studenten, die in Andreas Jungs Labor gearbeitet haben:

\paragraph{On simulations of thermal performance of the upgraded pixel detector}
My primary responsibilities include producing thermal simulations of anything the group needed, such as a replication of our carbon fiber calibration experiment or a more current design of our CO2 cooling structure.
Whenever Professor Jung begins an experiment involving heated detectors and fibers, he often relies on me to create (mostly from scratch) thermal simulations for comparison.
For the past year and a half, designs of CMS inner detectors by physicists at the LHC have been sent to us for simulation.
We perform the necessary tests and finite element analyses of their designs with some of Professor Jung’s suggestions and report our findings to the experts at Fermilab and CERN.

\paragraph{On work to built prototypes of pixel detector supports}
As an undergraduate researcher at Purdue Silicon Detector Labs I do a large variety of work, from programming to machining.
On a day to day basis my work generally consists of the measurement of carbon fiber laminate dimensions on a coordinate measuring machine.
Measurement include the thickness, done with a touch probe, and bow, done with an optical probe.
The process of measuring these values was originally largely manual and very time consuming.
One of my tasks is to make the process more efficient and automated.
I also am the labs gantry operator.
This means that I operate and program the robotic gantry used for a variety of precision applications.
I also participate in the production of the carbon fiber sheets from cutting to layout to curing.

Im folgenden erhältst du eine Liste mit typischen Aufgaben, die für dich geeignet sein könnten.
So kannst du dir schon mal ein Bild davon machen, was dich erwartet und vielleicht findest du auch schon ein Thema welches dich interessiert.
\begin{itemize}
  \item Visual Inspections of HDI/Modules (mostly completed since phase I
  upgrade is completed)
  \item IV tests of modules
  \item Gluing/encapsulation on the gantry
  \item Coding for gantry and other tooling (Labview)
  \item Coding for data analysis (ROOT, C++)
  \item Using CMM to measure parts/features (sensors and CF)
  \item Measure out/Lay up/ Cure CF plys
  \item Test thermal conductivity of Carbon foam and fiber
  \item Heat transfer of thermal wing (experiment)
  \item Finite Element Analysis on many structures (ANSYS)
  \item Design new tooling in CAD programs (Inventor, Solidworks, ANSYS)
  \item Machine new tools in the machine shop
  \item Soldering of electrical components
  \item Setting up new machinery/equipment
  \item Learn how to do experiments and work in a clean room
  \item Work with other Purdue research centers, e.g. CMSC and Birck
  \item Collaborate with Cornell/ other worldwide universities and institutions
  \item Learn how to work in a real lab environment and get experience without
  tremendous pressures
\end{itemize}

\subsubsection*{Erfahrungsberichte}
Erfahrungsberichte von ehmaligen Teilnehmern findest du auf unserer \href{https://pep-dortmund.org/stipendien/erfahrungsberichte.html} {Internetseite}.
\end{document}
